\documentclass{article}
\usepackage[spanish,es-noquoting]{babel}
\usepackage[utf8]{inputenc}
\usepackage{amsmath}
\usepackage{graphicx}
\usepackage{booktabs}
\usepackage{array}
\usepackage{circuitikz}

\title{Teorema de Thévenin - Práctica 4}
\author{}
\date{}

\begin{document}

\maketitle

\section*{Introducción}
a un circuito con una o más fuentes de alimentación.

6. Una vez reemplazada la red compleja por el circuito equivalente de Thévenin, la solución que da la corriente en la carga se puede obtener por la ley de Ohm. Así en la figura 26-1 d la corriente en la carga de 350 $\Omega$ es igual a $V_{TH}$ dividida por la suma de $R_{TH}$ y $R_L$:

\[
I_{\text{carga}} = \frac{V_{TH}}{R_{TH} + R_L}
\]

7. En esta Práctica verificaremos experimentalmente el teorema de Thévenin para un juego de valores de un circuito puente desequilibrado.

\section{AUTOEXAMEN}

\begin{enumerate}
    \item Se considera el circuito de la figura 26-1 a. $V = 160 \, V$, $R_1 = 30 \, \Omega$, $R_2 = 270 \, \Omega$, $R_3 = 500 \, \Omega$, $R_L = 560 \, \Omega$. En el circuito equivalente de Thévenin se supone que la resistencia interna de la fuente es cero.
    \begin{enumerate}
        \item $V_{TH} = \ldots \ldots \ldots \, V$.
        \item $R_{TH} = \ldots \ldots \ldots \, \Omega$.
        \item $I_{\text{carga}} = \ldots \ldots \ldots \, A$.
    \end{enumerate}
    
    \item Se considera el circuito de la figura 26-2 a. $V = 120 \, V$, $R_1 = 200 \, \Omega$, $R_2 = 500 \, \Omega$, $R_3 = 300 \, \Omega$, $R_4 = 600 \, \Omega$ y $R_s = 112,5 \, \Omega$. En el circuito de Thévenin equivalente, suponiendo que la resistencia interna de la fuente es cero,
    \begin{enumerate}
        \item $V_{TH} = \ldots \ldots \ldots \, V$.
        \item $R_{HT} = \ldots \ldots \ldots \, \Omega$.
        \item $I_s = \ldots \ldots \ldots \, A$, donde $I_s$ es la corriente en $R_s$.
    \end{enumerate}
\end{enumerate}

\begin{figure}[h!]
    \centering
    \begin{circuitikz}[scale=0.9, transform shape]
        % Definir coordenadas de los nodos del puente
        \coordinate (A) at (0, 2.5);
        \coordinate (B) at (-3, 0);
        \coordinate (C) at (3, 0);
        \coordinate (D) at (0, -2.5);

        % Dibujar los resistores del puente
        \draw (A) to[R, l_=$R_1$, l^=$3.300\,\Omega$] (B);
        \draw (A) to[R, l_=$R_2$, l^=$15.000\,\Omega$] (C);
        \draw (B) to[R, l_=$R_4$, l^=$5.600\,\Omega$] (D);
        \draw (C) to[R, l_=$R_3$, l^=$10.000\,\Omega$] (D);

        % Dibujar la rama central (carga)
        \draw (B) to[R, l_=$R_L$, l^=$2.700\,\Omega$] ++(2.5,0)
                  to[closing switch, l=$S_2$] ++(1.5,0)
                  to[ammeter, l=M] (C);

        % Dibujar la fuente de alimentación y el interruptor S1
        \draw (D) to[closing switch, l_=$S_1$] ++(-2,0)
                  to[short] ++(-1,0) coordinate (V_neg)
                  (A) to[short] (A |- V_neg)
                  to[V, v<={$V=40V$}] (V_neg);

        % Etiquetas de los nodos
        \node[above] at (A) {A}; \node[left] at (B) {B};
        \node[right] at (C) {C}; \node[below] at (D) {D};
    \end{circuitikz}
    \caption{Circuito experimental para confirmar el teorema de Thévenin.}
    \label{fig:26-3}
\end{figure}

\section{MATERIALES NECESARIOS}

\textbf{Equipo:} Fuente de c.c. variable regulada; EVM; miliamperímetro (o VOM) con un alcance de 5 a 10 mA. \\
\textbf{Resistencias:} $\frac{1}{2}$ W 2.700, 3.300, 5.600, 10.000 y 15.000 $\Omega$; y otras que requiera la operación 7 (discrecional). \\
\textbf{Diversos:} Caja de décadas de resistencia o, si no se dispone de ella, un potenciómetro de 10.000 $\Omega$ 2 W; dos interruptores unipolares.

\section{PROCEDIMIENTO}

\begin{enumerate}
    \item Conectar el circuito de la figura 26-3. $S_1$ y $S_2$ están abiertos. M es un miliamperímetro conmutado en el alcance de 1 mA. Cerrar $S_1$. Ajustar $V$ para que la tensión medida con un voltímetro electrónico sea 40 voltios. Cerrar $S_2$. Medir la corriente en $R_L$ y anotarla en la tabla 26-1 en la columna «$I_L$ medida circuito original».
    
    \item Abrir $S_2$. El interruptor $S_1$ permanece cerrado. Medir con un voltímetro electrónico la tensión $V_{BC}$ entre B y C y anotarla en la tabla 26-1 en la columna «$V_{TH}$ medida».
    
    \item Suprimir la fuente de tensión $V$ del circuito. $S_2$ está aún abierto. Cortocircuitar los puntos A y D.
    
    \textbf{NOTA:} Se supone que la resistencia interna de la fuente de tensión regulada es muy baja, y puede ser despreciada.
    
    $S_2$ está aún abierto. Medir la resistencia entre los puntos B y C y anotarla en «$R_{TH}$ medida».
    
    \item Ajustar la salida de la fuente de alimentación para $V_{TH}$ medida. Ajustar la resistencia de la caja de décadas al valor $R_{TH}$ y conectarla entre los puntos P y B. Si no se dispone de caja de resistencias, conectar una resistencia de 5.600 $\Omega$ en serie con un potenciómetro de 10.000 $\Omega$ que se utiliza como reóstato, y ajustar el reóstato para que la resistencia entre los puntos P y C (fig. 26-4) sea igual al valor medido de $R_{TH}$. Conectar ahora $R_L$, resistencia de carga de 2.700 $\Omega$, en serie con el miliamperímetro.
\end{enumerate}

% Tabla 26-1 (placeholder)
\begin{table}[h]
\centering
\caption{Tabla 26-1: Resultados experimentales}
\begin{tabular}{|c|c|c|}
\hline
$I_L$ medida circuito original & $V_{TH}$ medida & $R_{TH}$ medida \\
\hline
 & & \\
\hline
\end{tabular}
\end{table}

\end{document}